\Needspace{12\baselineskip}
\originalchapter{铷、无反冲谱和光阱公开预习题}
本章逐项翻译exp11的6题、exp13的6题、exp48的5题和exp51的3题。exp49的5题已在NMR量子信息章解答，合计25道实际公开预习题。其余手册所指但注明OCW用户不可用的网页问题没有恢复或编造。本章答案为编者推导；量级估算与仪器实测值明确分开。

\originalsection{exp11：光学泵浦的六个准备问题}
\begin{exercise}\label{task:exp11-1}
参照铷能级，估算以下差的eV数量级：(a)$5^2S_{1/2}$与 $5^2P_{3/2}$；(b)$5^2P_{1/2}$与 $5^2P_{3/2}$；(c)$5^2S_{1/2}$的F为2与3；(d)1 gauss场中F为2的 $m_F=-1$ 与0。
\end{exercise}
\begin{solution}
F为2、3指核自旋 $I=5/2$ 的 $^{85}$Rb。D2约780nm，光学差 $hc/\lambda\simeq1.59$eV。D1约794.8nm，所以细结构差 $hc(1/780-1/794.8)\,\mathrm{nm}^{-1}\simeq.0296$eV。基态超精细分裂约3.036GHz，对应 $1.26\times10^{-5}$eV；F为2的 $|g_F|\simeq1/3$，相邻Zeeman差 $|g_F|\mu_B10^{-4}\mathrm T\simeq1.93\times10^{-9}$eV。四级尺度相差极大，能级示意图不能用同一比例。
\end{solution}
\begin{exercise}\label{task:exp11-2}
用角动量与磁矩相加的向量模型或矩阵方法，推导两个铷同位素基态Landé因子。
\end{exercise}
\begin{solution}
基态L为0、$J=S=1/2$、$\mathbf F=\mathbf I+\mathbf J$。在固定F、m态上只有沿F的平均投影存留。利用 $2\mathbf J\cdot\mathbf F=F^2+J^2-I^2$，定义
\[
 c_J=\frac{F(F+1)+J(J+1)-I(I+1)}{2F(F+1)},\qquad c_I=1-c_J.
\]
取Zeeman Hamilton量 $H_Z=\mu_BB(g_JJ_z-g_n\mu_N I_z/\mu_B)/\hbar$，则 $g_F=g_Jc_J-g_n(\mu_N/\mu_B)c_I$。忽略小核项、$g_J\simeq2$，上支 $F=I+1/2$ 给 $g_F=1/(I+1/2)$，下支给负号。因此85Rb F为3、2分别约 $+1/3,-1/3$；87Rb I为3/2，F为2、1分别约 $+1/2,-1/2$。谱测正频率给绝对值，不能据此认为符号都正。若保留核项须核对磁矩归一化：手册写的1.3527、2.7505接近以核磁子表示的伸展态核磁矩，按本式的定义须再除I，不能直接当 $g_n$。独立参数参考以Bohr磁子计的 $g_I=-.0002936400,-.0009951414$，则直接用 $g_F=g_Jc_J+g_Ic_I$。
\end{solution}
\begin{exercise}\label{task:exp11-3}
320K热平衡下，一摩尔Rb有多少原子在 $5^2P_{1/2}$？1 gauss场中基态最低与最高磁子态人口差多少？
\end{exercise}
\begin{solution}
只取电子S、P1/2两个具有相同电子简并度2的多重态，核简并相同，激发总人口约
\[
 N_A\exp[-hc/(794.8\mathrm{nm}\,k_BT)]\simeq.163\quad\text{个每摩尔}.
\]
这是热平衡期望，甚至小于一个，不应当作每摩尔固定有小数原子；光泵下的人口远非该值。

磁子态的“最高/最低”须指定同一个F多重态；若把F为2、3的超精细差也算入，就不再是单纯磁分裂。85Rb基态共12子态，高温下每个约 $N_A/12$。F为3的极端m为 $\pm3$，差 $2\mu_BB$，所以
\[
 N_{\rm low}-N_{\rm high}\simeq\frac{N_A}{12}\frac{2\mu_BB}{k_BT}
 \simeq2.11\times10^{16}.
\]
相对全摩尔仅约 $3.5\times10^{-8}$，虽绝对数大，极化却极小。F为2相同计算的磁差为 $4\mu_BB/3$，人口差约 $1.40\times10^{16}$。87Rb用8基态子态及其g因子重新计算；天然同位素混合需再按丰度加权，不能把85Rb每摩尔结果冒充所有组成。
\end{solution}
\begin{exercise}\label{task:exp11-4}
给两个轴未知的线偏振片与一个轴未知的四分之一波片，怎样产生圆偏振光？
\end{exercise}
\begin{solution}
先让两偏振片交叉消光，确定其相对轴，任选第一片轴为0。将波片置于两片之间并旋转，消光位置对应其快/慢轴与第一片轴平行；旋转45度后两正交等幅分量延迟 $\pi/2$，出射为圆偏振。撤下第二片作输出；再旋转作为分析器时透射强度应近恒，不等于无偏振光，因为插回另一波片可还原线偏振并消光。该操作仅确定圆偏振而不唯一确定左右旋，需知道波片快轴及传播方向；真实波片的设计波长、偏角和损耗会使输出椭圆化。
\end{solution}
\begin{exercise}\label{task:exp11-5}
Rb核对(a)总角动量、(b)总磁矩的贡献各占多少？
\end{exercise}
\begin{solution}
“比例”取沿F的投影而非长度相除。上题 $c_I$ 给85Rb上/下支 $5/6,7/6$，87Rb上/下支 $3/4,5/4$；下支核投影超过1，是因为电子投影反向，非算错。磁矩比例则为 $[-g_n(\mu_N/\mu_B)c_I]/g_F$，核磁矩虽角动量占主要部分，却因 $\mu_N/\mu_B\simeq1/1836$ 通常只在千分之几；上支与电子项反号、下支可同号。若只说“核贡献很小”，会混淆角动量与磁矩这两件事。
\end{solution}
\begin{exercise}\label{task:exp11-6}
把地磁当作地心、沿自转轴的理想点偶极，Boston纬度处场指向哪里？
\end{exercise}
\begin{solution}
球坐标偶极场 $B_r\propto2\cos\theta,B_\theta\propto\sin\theta$，纬度 $\lambda=\pi/2-\theta$ 给相对水平倾角 $\tan I=2\tan\lambda$。取Boston约42.4度，$I\simeq61.3$度。按地球北半球场向北且向下的实际极性，指向磁北并下倾；若假设偶极矩反向，全部箭头反向。理想轴对称模型不给实际磁偏角，不能用它替代当地实测地磁。
\end{solution}

\originalsection{exp13：Mössbauer的六个准备问题}
\begin{exercise}\label{task:exp13-1}
质量M的核发出能量E的gamma光子，推导准确相对论反冲能；铁核、E为14.4keV时求值。
\end{exercise}
\begin{solution}
在初始静止系，反冲动量大小 $E/c$。若M指末态核静质量，$T=\sqrt{M^2c^4+E^2}-Mc^2$；若M指初态质量，能量守恒给 $T=E^2/[2(Mc^2-E)]$。两种定义不能混用；在 $E\ll Mc^2$ 都约 $E^2/(2Mc^2)$。用57原子质量单位估核质量，约 $1.95\times10^{-3}$eV，比自然线宽约 $4.66\times10^{-9}$eV大约四十万倍。自由核发射与吸收分别反向偏移，因此共振失配约两个反冲能。
\end{solution}
\begin{exercise}\label{task:exp13-2}
同一个核嵌于自由移动、质量 $10^{-6}$g的铁小晶体，求整体反冲能。无相对晶体质心的核反冲时，gamma光子因整体反冲改变多少能量？
\end{exercise}
\begin{solution}
晶体质量 $10^{-9}$kg；以总质量代替单核，反冲能约 $1.85\times10^{-19}$eV。光子比内部跃迁能少此值，远小于自然宽。“无反冲”指无内部声子激发、没有可辨的单核反冲能损失，不是违反总动量守恒。只部分事件满足；无反冲比例由晶格振动与温度决定。
\end{solution}
\begin{exercise}\label{task:exp13-3}
画57Fe核在磁场中能级图，并标出本实验Zeeman谱允许跃迁。
\end{exercise}
\begin{solution}
基态 $I_g=1/2$ 两子态，激发态 $I_e=3/2$ 四子态，分别 $E_i=-g_i\mu_NBm_i$。磁偶极跃迁 $\Delta m=0,\pm1$：$3/2\to1/2,1/2\to1/2,1/2\to-1/2,-1/2\to1/2,-1/2\to-1/2,-3/2\to-1/2$，共六条。能量 $E_0-g_e\mu_NBm_e+g_g\mu_NBm_g$。下图仅示连接，不按实际g比、正负或线能顺序绘制。
\begin{center}\begin{tikzpicture}[x=1cm,y=.75cm]
\draw (0,4)--(1.6,4) node[right]{$m_e=3/2$};\draw (0,3)--(1.6,3) node[right]{$1/2$};
\draw (0,2)--(1.6,2) node[right]{$-1/2$};\draw (0,1)--(1.6,1) node[right]{$-3/2$};
\draw (5,0)--(6.6,0) node[right]{$m_g=1/2$};\draw (5,-1)--(6.6,-1) node[right]{$-1/2$};
\draw[->,main] (1,4)--(5.3,0);\draw[->,main] (1,3)--(5.5,0);\draw[->,blue] (1,3)--(5.5,-1);
\draw[->,main] (1,2)--(5.7,0);\draw[->,blue] (1,2)--(5.7,-1);\draw[->,blue] (1,1)--(5.9,-1);
\end{tikzpicture}\end{center}
粉末无偏振理想强度常为3:2:1:1:2:3；取向、极化和吸收厚度会改变强度，不能强行固定所有实谱比。
\end{solution}
\begin{exercise}\label{task:exp13-4}
为什么单独电四极场梯度的核能级重数小于磁偶极Zeeman分裂？
\end{exercise}
\begin{solution}
轴对称四极能 $E_Q\propto3m^2-I(I+1)$，对m正负相同，而磁能正比m。故激发态3/2只有 $|m|=1/2,3/2$ 两组，基态1/2没有四极矩，典型给双线；磁场一般将全部m拆开。非轴对称项混合m，半整数自旋在无磁、保持时间反演下仍有Kramers二重简并。磁场破坏这种保护，所以不能把四极看成较弱版同样的四级均分。
\end{solution}
\begin{exercise}\label{task:exp13-5}
源相对吸收体需多快，才把共振频率移一个自然线宽？
\end{exercise}
\begin{solution}
手册 $t_{1/2}=9.8\times10^{-8}$s。平均寿命 $\tau=t_{1/2}/\ln2$，$\Gamma=\hbar/\tau=\hbar\ln2/t_{1/2}\simeq4.66\times10^{-9}$eV；因此 $v=c\Gamma/E\simeq9.69\times10^{-5}$m/s，约.097mm/s。原手册文字、旧文献的lifetime标法不一致，此处明确半衰期与平均寿命。实际源和吸收Lorentz宽卷积会相加，有限厚度再展宽，不应将观测卷积宽直接当单核自然宽。
\end{solution}
\begin{exercise}\label{task:exp13-6}
为什么不能用源同时发出的6.4keV光子做同一Mössbauer实验？
\end{exercise}
\begin{solution}
6.4keV来自原子壳层填空的特征X射线，而14.4keV是长寿命核态间跃迁。前者的短空穴寿命给宽原子谱线，而且没有这对57Fe长寿命核态的无反冲共振吸收。能量更低并不自动使它成为可用的核Mössbauer共振。
\end{solution}

\originalsection{exp48：饱和吸收的五个准备问题}
\begin{exercise}\label{task:exp48-1}
波长λ附近小变化 $\Delta\lambda$ 对应怎样的 $\Delta\nu$？λ为780nm时用GHz/nm给换算。
\end{exercise}
\begin{solution}
$\nu=c/\lambda$，$\Delta\nu=-c\Delta\lambda/\lambda^2$，780nm处斜率约 $-492.755$GHz/nm。负号表示波长增加时频率降低；大范围应使用完整倒数式，不能无限延伸局部斜率。
\end{solution}
\begin{exercise}\label{task:exp48-2}
长度L、折射率n的激光腔，λ处纵模频率与波长间距是多少？分别算780nm、(a)L为250µm且n为3.5的二极管；(b)L为2cm且n为1的外腔。
\end{exercise}
\begin{solution}
往返相位 $2nL\nu/c$ 为整数，忽略色散给 $\Delta\nu=c/(2nL)$，小间距的波长大小 $|\Delta\lambda|=\lambda^2/(2nL)$。两者约171.310GHz/.347657nm和7.49481GHz/.015210nm。真实色散以群折射率决定模间距，外腔反馈与二极管内腔同时存在，不能仅凭外腔长度断言任意工作点都是单模。
\end{solution}
\begin{exercise}\label{task:exp48-3}
画85Rb、87Rb基态5S与首激发5P的精细、超精细结构，估各相互作用尺度。
\end{exercise}
\begin{solution}
$L=0$ 的5S仅 $J=1/2$；5P有 $J=1/2,3/2$，精细差约.03eV，即约7THz。每支 $F=|I-J|,\ldots,I+J$。下图自绘结构，所有竖向间隔不按比例。
\begin{center}\begin{tikzpicture}[x=1cm,y=.65cm]
\node at (1,6.2){85Rb，$I=5/2$};\node at (5.2,6.2){87Rb，$I=3/2$};
\foreach \xx in {0,4.2}{\draw (\xx,0)--++(1.8,0);\draw (\xx,.6)--++(1.8,0);\draw (\xx,3)--++(1.8,0);\draw (\xx,3.6)--++(1.8,0);\draw (\xx,4.6)--++(1.8,0);\draw (\xx,4.9)--++(1.8,0);\draw (\xx,5.2)--++(1.8,0);\draw (\xx,5.5)--++(1.8,0);}
\node[left] at (0,.3){5S};\node[left] at (0,3.3){5P$_{1/2}$};\node[left] at (0,5.05){5P$_{3/2}$};
\node[right] at (1.8,0){$F=2$};\node[right] at (1.8,.6){3};\node[right] at (1.8,3){2};\node[right] at (1.8,3.6){3};
\node[right] at (1.8,4.6){1};\node[right] at (1.8,4.9){2};\node[right] at (1.8,5.2){3};\node[right] at (1.8,5.5){4};
\node[right] at (6,0){$F=1$};\node[right] at (6,.6){2};\node[right] at (6,3){1};\node[right] at (6,3.6){2};
\node[right] at (6,4.6){0};\node[right] at (6,4.9){1};\node[right] at (6,5.2){2};\node[right] at (6,5.5){3};
\end{tikzpicture}\end{center}
基态分裂85Rb约3036MHz、87Rb约6835MHz；5P1/2约362、817MHz；5P3/2相邻间隔约几十至几百MHz。超精细量级不是所有F间隔相等；精确能量由磁偶极A和四极B项给出，须用各支的参数，不把仅示结构的图当数值表。估算参数与Daniel Steck的公开D-line数据核对，仅引事实与链接，未复制其表或图。
\end{solution}
\begin{exercise}\label{task:exp48-4}
以MHz列表比较激光线宽、纵模间距、原子自然宽、精细和超精细差、Doppler宽、Fabry--Perot自由谱区与分辨率、波长计分辨率等尺度。
\end{exercise}
\begin{solution}
按原题条件及前两题的换算，各尺度汇总如下。
\begin{center}\begin{tabular}{p{.49\linewidth}|p{.4\linewidth}}量与条件&近似尺度（MHz）\\\hline
手册自由运行激光宽&30（实际Lamb凹陷约30--40）\\
二极管内腔/上述2cm外腔间距&171310 / 7495\\
Rb D2自然FWHM&6\\
5P精细差&$7\times10^6$\\
基态超精细85/87Rb&3036 / 6835\\
297K、780nm Doppler FWHM&手册约502\\
48.000cm参考腔FSR&手册参考曲线312.197\\
Fabry--Perot分辨率&FSR/实测精细度 $\mathcal F$\\
波长计约±.2nm&约±98551
\end{tabular}\end{center}
参考腔312.197是原手册研究级参考图条件，实际实验须量自己腔长；原题没有给自己腔的精细度，不能伪造其分辨率，可由单峰宽/相邻间距求。波长计粗读数不能标定MHz级超精细间隔。Doppler宽公式 $\Delta\nu_D=(\nu_0/c)\sqrt{8k_BT\ln2/M}$；光强加大可提高信号也会功率展宽。
\end{solution}
\begin{exercise}\label{task:exp48-5}
解释二极管激光的增益介质、泵能来源及谐振光腔由什么组成。
\end{exercise}
\begin{solution}
正向电流向半导体有源区注入电子/空穴，非平衡载流子使光学跃迁具有净受激发射增益。芯片两端面反射提供内腔往返反馈，增益大于损耗才达到阈值；不是简单把任意LED当作单色激光。外腔方案另用准直与光栅选择反馈波长，光栅、二极管端面和外腔构成复合反馈。电流与温度同时改变载流子增益和光程，压电调整外腔光程；不同作用可导致模跳，连续扫描需要记录单模范围。
\end{solution}

\originalsection{exp51：光阱的三个准备问题}
\begin{exercise}\label{task:exp51-1}
在光线光学极限用折射光动量变化解释稳定陷阱的横向、纵向恢复力；定性画轻微位移后的力。阱中心是否与光焦点相同？为何用高NA？给定本实验激光与粒径，光线模拟可定量相信吗？原题另建议探索外部Java模拟器。
\end{exercise}
\begin{solution}
折射使射线向球内偏折，球受相反动量变化。横向移出强度轴时，近轴侧较强光贡献恢复力，$F_x\simeq-k_xx$；纵向焦区的梯度力与向传播方向的散射推力相加，平衡点通常在几何焦点下游，近平衡 $F_z\simeq-k_z(z-z_0)$。高NA使射线角度大、轴向梯度强，才能抵消散射而形成三维稳定阱。
\begin{center}\begin{tikzpicture}[x=1.2cm,y=1cm]
\draw[->] (-2,0)--(2,0) node[right]{位移};\draw[->] (0,-1.3)--(0,1.3) node[above]{恢复力};\draw[main,thick] (-1.4,1)--(1.4,-1);\node at (1.2,1){平衡点附近};
\end{tikzpicture}\end{center}
手册粒径约3.2µm、波长近1µm，粒径仅数倍波长，几何光线可解释方向但不保证定量力或各方向刚度。定量需波动电磁散射与介质参数、实际焦场；本项目未运行已不可在当前浏览器直接使用的旧Java小程序，也未把定性图冒充其输出。
\end{solution}
\begin{exercise}\label{task:exp51-2}
估水中典型细菌停止的时间和距离，比较生物尺度；平均减速度的ma与维持初速的力怎样比较？
\end{exercise}
\begin{solution}
明确作球形量级模型：半径1µm、密度 $10^3$kg/m$^3$、水黏度 $10^{-3}$Pa s、初速20µm/s。质量 $m=4\pi\rho a^3/3$，阻力 $\zeta=6\pi\eta a$，停驱动后 $v=v_0\e^{-t/\tau_m}$，$\tau_m=m/\zeta=2\rho a^2/(9\eta)\simeq2.22\times10^{-7}$s，总惯性滑行距离 $v_0\tau_m\simeq4.44\times10^{-12}$m。指数严格不在有限t变0，“停止时间”须给阈值，如降至1\%约 $4.61\tau_m$。这些远小于细菌移动一体长的约.1s和微米长度。

初始减速度 $a_0=-v_0/\tau_m$，所以 $m|a_0|=\zeta v_0\simeq.377$pN，等于维持该速度所需的阻力。以一惯性时间的平均减速则小一个 $1-\e^{-1}$ 因子；一般平均ma取决于时间窗，不能将它写成任意定义下完全相等。实际细菌细长体、近壁和流体附加质量可改量级系数，低Re结论不靠球模型精确成立。
\end{solution}
\begin{exercise}\label{task:exp51-3}
本实验主要危险是什么，应怎样避免？
\end{exercise}
\begin{solution}
按原手册，红外激光不可见，主风险为眼损伤；使用匹配波长/光密度的护目镜、封闭或挡住束及反射、按实验室程序调光，不用眼观察激光束。制样的碎盖玻片、刀具属生物锐器，手套、合适废弃容器和监督下制样；细菌等样品按设施的生物规程处理。电子设备和显微镜运动部件按原装置指导操作。本答案是公开预习题的功能说明，不能替代当地培训与装置操作许可。
\end{solution}
